% ── encoding.tex · §2 ──────────────────
\section{Encoding \pp{} in \lp}\label{sec:encoding}

\subsection{The Predicate Prover}\label{sec:pp-syntax}

Syntax for the \emph{predicates} and \emph{expressions} of \pp{} are
given as follows:
%
% do we really have n-ary function applications?
% do we really have n-ary membership?
% is this the full grammar for predicates/expressions?
\begin{align*}
  E, F &\Coloneqq x \mid n \mid E \mapsto F \mid f(E_1,\ldots,E_k)
         \mid E + F \mid E - F \mid -E \\
  P, Q &\Coloneqq E \mid E = F \mid E \leq F
         \mid E_1,\ldots,E_k \in F
         \mid \lnot P \mid P \land Q \mid P \lor Q \mid P \Rightarrow Q
         \mid P \Leftrightarrow Q
         \mid \mathcal{Q}\,\vec{x} \cdot P
\end{align*}
where $x$ ranges over identifiers, $n$ over numerals, $\mapsto$ is B's
pairing, and $\mathcal{Q} \in \{\forall, \exists, \Dia, \Hrt\}$ binds a
list of variables.
%
The quantifiers $\Dia$ and $\Hrt$ are universal quantifiers idiosyncratic
to \pp's normalisation procedure; semantically both are $\forall$ and differ
only in the role they play during proof search.
%
% nonsense sentence? maybe its worth explaining the pre-processing step.
% although its clear that some set-theory stuff remain as function/relation
% symbols, but the rules don't deal with them in any meaningful way.
The set-theoretic surface of B (powersets, intervals, relational
operators, \ldots) reaches \pp{} as applications of function symbols;
\pp{} reasons about them through its normalisation and equality rules,
not through their set-theoretic definitions.

% we need to formally introduce sequents and results here before
% discussing the rules.
\pp's inference rules are specified in the annex of its internal
specification~\cite{pp_spec}: roughly \numRulesSpec{} rules, grouped by
families — propositional (\rn{AND}, \rn{OR}, \rn{IMP}, \rn{EQV},
\rn{NOT}), hypothesis discharge (\rn{AXM1--9}), quantifiers
(\rn{ALL1--9}, \rn{XST1--8}), equality and the equality prover
(\rn{EQL}, \rn{EAXM}, \rn{EQC}, \rn{ECTR}, \ldots), booleans
(\rn{BOOL}), linear arithmetic (\rn{AR1--13} plus the solver terminals
\rn{ARITH} and \rn{EGALITE}), and normalisation and instantiation
(\rn{NRM1--30}, \rn{INS}).
%
Each rule application reduces the current goal to zero, one,
or two premises, possibly under computational side-conditions
(``$P$ occurs in the list of hypotheses $H$'',
 ``$a$ and $b$ are integer literals such that $a + b < 0$'', etc.).

%  everything here is kind of filler and has already been said.
% we should show an example replay (involving ALL7) and how it
% corresponds to an 'abstract derivation'.
How \pp{} \emph{finds} a proof is irrelevant here: we consume its
replays.
%
\pp{} can be instrumented to record a compact \trace{} of the rules it
applied; a companion utility (\textsc{replay}) re-runs the proof and
expands the trace into a \replay{} that annotates every step with the
formula it was applied to:
\begin{codePP}
[RULE] <formula>        rule application + the goal it rewrote or closed
[RULE(a)] <formula>     rule argument (hypothesis, subscript, predicate)
[RULE_1] <formula>      result-form variant, inside a normalisation chain
[RULE_n] <formula>      n-ary variant binding n variables
[FIN(P)], [STOP_NORM]   phantom normalisation markers
\end{codePP}
The main proof is recorded prefix-style (a rule precedes its
subproofs), with two irregularities: the \emph{result chain} feeding a
branching quantifier (\rn{ALL7}, \rn{XST8}) is emitted \emph{before}
the branch rule itself, and phantom normalisation markers are
interleaved throughout.
%
% the notion of 'phantom' is useless jargon. remove.
The first non-phantom annotation is the overall goal.

% result should be explained earlier.
The \texttt{\_1} (\emph{result-form}) variants deserve explanation, as
they shape the encoding.
% i prefer the framing that results/sequents are just different
Inside a normalisation chain a rule does not close a sequent: it
transforms a formula and passes the \emph{result} on.
%
The specification (\S 8.13) assigns every result-form rule one of three
schemas: its result is the constant true (the formula is discharged), a
single transformed formula, or the combination of the results of its two
children.

\subsection{\lp}\label{sec:lp}
% we should have a grammar to mirror what we showed for PP.
% for the lambdapi commands, we should use the same syntax highlighting we used in eo2lp
\lp~\cite{hondet:hal-02981561} is a proof assistant for the
\lpcmr~\cite{CousineauDowek}, the calculus underlying the \dedukti{}
logical framework~\cite{assaf2023dedukti}: a dependently-typed
$\lambda$-calculus whose conversion is generated by $\beta$ together
with user-declared rewrite rules.
%
A development is a sequence of declarations.
%
\texttt{symbol} introduces a typed constant, with explicit
$(x : A)$ and implicit $[x : A]$ parameters and an optional definition;
\texttt{rule} $\ell ↪ r$ adds a rewrite rule over pattern variables
\texttt{\$x}, checked for subject reduction at declaration time;
\texttt{inductive} generates constructors and a recursor.
%
Modifiers refine a symbol's status: \texttt{constant} symbols admit no
rewrite rule at their head, \texttt{injective} symbols may be inverted
during unification, and the definition of an \texttt{opaque} symbol is
sealed once type-checked.
%
A definition may be given as a proof script (\texttt{begin} \ldots{}
\texttt{end}) whose tactics --- \texttt{assume}, \texttt{refine},
\texttt{rewrite}, \texttt{have}, \texttt{induction} --- ultimately
elaborate to a proof term that the kernel re-checks, so the tactic layer
adds nothing to the trusted base.

We build on \lp's standard library~\cite{Blanqui2021}: object-level types are
encoded in a Tarski-style universe $\Setu$ with $\El : \Setu \to \mathsf{TYPE}$
mapping each set code to its type of inhabitants, and a type $\Prop$ of
propositions with $\Prf : \Prop \to \mathsf{TYPE}$ mapping each proposition
to its type of proofs (the library also provides the code $o : \Setu$ with
$\El\,o ↪ \Prop$).
%
Connectives and quantifiers are the standard impredicative encodings;
on top of these we assume the library's classical axioms — excluded
middle and propositional extensionality.

\subsection{Encoding PP's Predicates and Expressions}\label{sec:enc-formulas}

The encoding is shallow.
% the explanation of the `o` set code is in the previous section when
% it naturally falls alongside the explanation of the `ι` set code.
All B expressions inhabit a single domain: a set code $\iota$, so an
expression is a term of type $\El\,\iota$ and a predicate is a term of
type $\Prop$.
%
The kernel file \texttt{B.lp} declares the domain constants —
$\mapsto$ (pairing), $\epsilon$ (membership), $\mathsf{BOOL}$ with
$\mathsf{BTRUE}$ and $\mathsf{BFALSE}$, B's functional application, and
the arithmetic signature $\mathbf{0}, \mathbf{1}, +, -, \leq$ — together
with their axioms: injectivity and congruence of pairing, set
extensionality, the two-element axioms for $\mathsf{BOOL}$, and for
arithmetic the theory of a discretely ordered abelian group
($+$ commutative and associative with identity $\mathbf{0}$ and inverse
$-$; $\leq$ a partial order, invariant under translation; and the
discreteness axioms $\mathbf{0} \leq \mathbf{1}$,
$\lnot(\mathbf{1} \leq \mathbf{0})$ and
$\lnot(x \leq \mathbf{0}) \Rightarrow \mathbf{1} \leq x$).
%
This is just enough to \emph{derive} \pp's arithmetic rules; it is not a
formalisation of $\mathbb{Z}$.
%
Numerals are decimal literals of the library's $\mathbb{N}$, coerced
into the domain by an embedding that \lp{} inserts automatically.
%
Every other operator (set operators, projections, \ldots) is an
uninterpreted constant.

\paragraph{$n$-ary conjunction.}
\pp{} treats $\land$ as left-associative and its rules address a
conjunction through its \emph{last} conjunct and the remainder; worse,
when a normalisation step discharges a conjunct, \pp{} silently drops it
from every formula it subsequently records.
%
Encoding $\land$ binarily would therefore force proofs to rewrite
associativity and unit collapses at nearly every step.
%
Instead, a conjunction is a \emph{snoc list}: an inductive type
$\mathbb{L}^{*}$ with empty list $∎$ and right append $∷$, and an
injective operator $\bwand : \mathbb{L}^{*}\,o \to \Prop$ with the
rewrite rules
$\bwand(∎ ∷ P) ↪ P$ and
$\bwand(ps ∷ \top) ↪ \bwand ps$
(a singleton conjunction \emph{is} its conjunct; discharged conjuncts
erase).
%
Its structural interface — introduction, and elimination of the last
conjunct and of the remainder — is axiomatised.
%
\pp's conjunction rules then become single structural matches; for
example (\texttt{lp/rules/Conj.lp}):
\begin{codeLP}
opaque symbol AND2 [ps : $\mathbb{L}^{*}$ o] [P : Prop]
  : π ($\bigwedge$ ps ⇒ ¬ P) → π (¬ $\bigwedge$ (ps ∷ P)) ≔
begin
  assume ps P h1 hc;
  refine h1 ($\bigwedge$_init hc) ($\bigwedge$_last hc)
end;
\end{codeLP}

\paragraph{Compound binders.}
%  this is NOT true. PP doesn't regroup quantifiers.
% some rules do this regrouping but really we are just explaining that PP
% distinguishes between ∀ (x,y), P and ∀ x, ∀ y, P
\pp{} regroups consecutive quantifiers — $\forall x \cdot \forall y
\cdot P$ becomes $\forall (x,y) \cdot P$ — and its quantifier rules
manipulate the compound binder as a unit.
%
We mirror this with an inductive dependent type $\Tuple : ℕ → \TYPE$
with projections $\prj [n] : \Tuple\,n → \El\,\iota$, and quantifiers
$\allI, \xstI : (\Tuple\,n \to \Prop) \to \Prop$
whose introduction and elimination principles are axiomatized bridges to
the $\Pi$-types of \lp{} (e.g.\
$(\Pi v{:}\,\Tuple\,n,\ \Prf\,(P\,v)) \to \Prf\,(\allI v, P\,v)$ and its
converse); $\Dia$ and $\Hrt$ are defined aliases of $\allI$ with bridges
of their own.
%
% we probably don't need to make this point but it's useful to have:
% -------------------------------------------------------------------
% No rewrite rule unfolds $\allI$ to nested $\forall$: if the quantifier
% reduced away, the surface form would vanish from goals before any rule
% lemma could address it, and inferring the
% implicit predicate $P$ at rule applications would require inverting the
% reduction, a higher-order unification problem \lp{} does not attempt.
%

\subsection{Encoding Rules}\label{sec:enc-rules}

Each \pp{} rule is a \lp{} symbol whose type states the rule and whose
body \emph{proves} it.
%
The encoding of a rule's statement is positional: the conclusion sequent
$H \seq P$ appears as $\Prf\,P$; a premise that extends the hypotheses,
$H, P \seq Q$, appears as the function type $\Prf\,P \to \Prf\,Q$; and a
side-condition ``$P$ occurs in $H$'' becomes an explicit argument of
type $\Prf\,P$, which our translation tool supplies by looking the hypothesis
up in scope.
%
For example, \rn{IMP4} (``to prove $P \Rightarrow Q$, assume $P$ and
prove $Q$'') and the discharge rule \rn{AXM1} (``$\lnot P$ occurs in
$H$, so $P \Rightarrow Q$ holds'') have the following type signatures:
\[
  \rn{IMP4} : (\Prf\,P \to \Prf\,Q) \to \Prf\,(P \Rightarrow Q)
  \qquad
  \rn{AXM1} : \Prf\,(\lnot P) \to \Prf\,(P \Rightarrow Q),
\]


\paragraph{Result chains.}
To encode the result-form (\texttt{\_1}) rules we package a
normalisation result as a dependent record: for a formula $t$,
$\Res\,t$ pairs the list $ts$ of surviving conjuncts with a proof of
$\Prf\,(t = \mathsf{conj}\,ts)$, where $\mathsf{conj}$ folds a list
back into a formula.
%
The three specification schemas map directly onto its constructors: a
discharged formula is the empty list with a proof of $t = \top$; a
transformation reuses a proved equality law (for example, \rn{NOT2}'s
chain form is
$\mathsf{mk_1}\,(\restm\,r)\,(\mathsf{eq\_trans}\ \texttt{¬¬-law}\
(\mathsf{res\_eq}\,r))$); and a combination concatenates the two
children's lists, justified by a concatenation law for $\mathsf{conj}$.
%
The $\top$-erasure rewrite rules on $\bwand$ and on list concatenation
make the combined result reduce to exactly the conjuncts \pp{} records —
discharged conjuncts disappear in any position, mirroring \pp.
%
A branching quantifier consumes a chain.
%
\rn{ALL7} normalises the body of a universal hypothesis pointwise and
continues under the normalised form:
\begin{codeLP}
opaque symbol ALL7 [n] [P : Tuple n → Prop] [Q : Prop]
  (ρ : Π v : Tuple n, Res (P v))
  : π ((`$\diamondsuit$ v, res_tm (ρ v)) ⇒ Q) → π ((`!! v, P v) ⇒ Q) ≔ $\ldots$
\end{codeLP}
where the proof transports each instance along the chain's equality
($\mathsf{res\_eq}$).

\paragraph{Solver side-conditions.}
% we should actually give an example of a on-the-fly generated proof terms.
Some rules are licensed not by a structural premise but by a verdict of
\pp's arithmetic solver or equality prover: \rn{AR4} requires
$(E+F) > \mathbf{0}$ for terms \pp{} already believes positive, and the
\rn{ARITH} terminal closes a goal because the accumulated
$\leq$-hypotheses are contradictory.
%
The corresponding lemmas take these facts as explicit $\Prf$-arguments,
and \tool{} \emph{generates} the witnessing proof at emission time: a
nonnegative (Farkas-style) combination of the $\leq \mathbf{0}$
hypotheses refuting the \rn{ARITH} goal, equality-transport chains for
the equality prover, and permutation proofs from the group axioms
wherever \pp{} recorded a sum in a different association order than the
lemma needs.
%
There is no fallback oracle: when a witnessing proof cannot be
generated, emission fails loudly (\autoref{sec:trust}).

% the whole following subsection is probably not needed.
% we are not saying anything other than 'we proved everything'.
\subsection{The Trusted Base}\label{sec:trust}
A reconstructed proof type-checks against: the \lp{} kernel; the
standard library with its classical axioms; the B domain axioms of
\autoref{sec:enc-formulas}; the structural interface of $\bwand$; and
the quantifier bridges.
%
Everything else is a theorem.
%
Of the \numRulesCovered{} rules in the library, every rule that can
occur in a replay carries a soundness proof; the \numRulesTrust{}
remaining chain forms of the regrouping rules are still stated with a
postulated oracle (\texttt{trust}), but they are unreachable — the
translator regroups binders at the syntax level, so no replay path emits
them.
%
The emitter is incapable of producing a proof that mentions
\texttt{trust} at all: its term grammar has no such constructor, so any
side-condition it cannot prove is a hard failure rather than a silent
axiom.
